\section{总结与展望}

\subsection{实验总结}

本学期六组实验从基础环境配置开始，逐步完成机械臂运动控制、抓取搬运、相机标定和视觉引导抓取。整体技术路线可以概括为：

\[
\begin{aligned}
\text{环境配置}&\rightarrow \text{运动学建模}\rightarrow \text{轨迹规划}\\
&\rightarrow \text{抓取动作}\rightarrow \text{相机标定}
\rightarrow \text{视觉识别}\rightarrow \text{视觉抓取}.
\end{aligned}
\]

实验一完成了 Python、Robotics Toolbox、OpenCV、串口通信等基础环境配置；实验二建立 DH 模型并实现末端直线运动；实验三在示教点基础上完成抓取搬运；实验四通过棋盘格标定获得相机内参、畸变参数和工作平面坐标关系；实验五使用颜色分割完成目标定位和视觉抓取；实验六使用 YOLO 替代颜色阈值识别，提高了视觉前端的鲁棒性和可扩展性。

\subsection{误差来源分析}

综合各实验，系统误差主要来自以下几个方面：

\begin{itemize}
  \item \textbf{运动学模型误差：} DH 参数与真实机械臂结构之间存在制造和装配误差，舵机零位也可能与模型零位不完全一致。
  \item \textbf{执行机构误差：} 一体化关节存在回差、摩擦和负载扰动，夹爪闭合宽度也会受物体材质影响。
  \item \textbf{相机标定误差：} 棋盘格角点检测、相机固定稳定性、工作平面高度假设都会影响像素到世界坐标的转换。
  \item \textbf{视觉检测误差：} 颜色阈值受光照影响明显，YOLO 检测则依赖训练数据覆盖范围和标注质量。
  \item \textbf{任务规划误差：} 抓取高度、下降距离、目标角度估计和放置点高度若设置不合理，会导致抓取不稳或放置偏差。
\end{itemize}

\subsection{改进方向}

后续可以从以下方向继续完善系统：

\begin{enumerate}
  \item 增加手眼标定流程，减少相机坐标到机械臂基坐标之间的人工标定误差。
  \item 在抓取前加入二次视觉确认，检测物体是否在夹爪正下方。
  \item 引入深度相机或多视角测量，避免单平面高度假设带来的误差。
  \item 扩充 YOLO 数据集，加入不同光照、遮挡、背景和物体姿态。
  \item 将轨迹规划从关键点插值升级为带速度、加速度约束的更平滑规划。
  \item 在真实执行中记录关节反馈和末端误差，形成闭环校正机制。
\end{enumerate}

总体来看，本实验已经形成了一个可运行、可扩展的桌面机械臂视觉抓取系统。它的核心价值不只是完成单个抓取动作，而是建立了从感知、坐标转换、运动规划到执行控制的完整链路，为后续更复杂的自主操作任务打下基础。
